\chapter{Sigmoidoscopy}

\section*{Overview}
Sigmoidoscopy uses a flexible scope to examine only the rectum and the lower part of the colon (sigmoid colon), which constitutes the distal third of the large intestine.

\section*{Alternative Names}
Flexible sigmoidoscopy

\section*{Principle}
A flexible endoscope is passed through the anus into the rectum and sigmoid colon, which is inflated to allow inspection of the lining. Biopsies and treatment of lesions in this distal segment can be performed through the scope.
\section*{History}
Rigid sigmoidoscopes were used in the early 20th century. Flexible fiberoptic sigmoidoscopes were introduced in the 1970s, providing a more comfortable screening option.

\section*{Why This Test or Procedure Is Ordered}
Ordered for colorectal cancer screening, evaluating rectal bleeding, monitoring inflammatory bowel disease, or investigating changes in bowel habits.

\section*{Is This a Primary or Secondary Test or Procedure}
Alternative screening option for colorectal cancer, often combined with annual fecal tests.

\section*{How This Test or Procedure Is Performed}
You lie on your left side. Sedation is usually not required. The scope is inserted and advanced up to 60 cm into the lower colon.

\section*{What Preparation Is Needed}
Cleanse the lower bowel using one or two enemas on the morning of the procedure. A clear liquid diet may be recommended.

\section*{Contraindications and Precautions}
Acute severe diverticulitis, suspected colonic perforation, or toxic megacolon.

\section*{Risks}
Extremely low risk of perforation (<1 in 10,000) or bleeding (<1 in 2,000 if biopsy is performed).

\section*{What Happens After the Test or Procedure}
You can immediately drive and return to normal activities. You may pass gas to relieve mild bloating.

\section*{Understanding Your Results}
Allows visualization of hemorrhoids, fissures, polyps, or inflammation in the distal colon. Biopsies can be performed.

\section*{Normal Results}
Normal mucosal lining of the rectum and sigmoid colon; no polyps, inflammation, or lesions.

\section*{Follow-up Tests and Procedures}
Repeat sigmoidoscopy in 5 years, or full colonoscopy if polyps or lesions are identified.

\section*{References}
\begin{enumerate}
  \item MedlinePlus. Sigmoidoscopy. U.S. National Library of Medicine; 2025.
  \item Current Medical Diagnosis and Treatment. McGraw-Hill; 2025.
\end{enumerate}

\clearpage
